% $Header:
% /home/vedranm/bitbucket/beamer/solutions/conference-talks/conference-ornate-20min.en.tex,v
% 90e850259b8b 2007/01/28 20:48:30 tantau $

\documentclass{beamer}

% This file is a solution template for:

% - Talk at a conference/colloquium.
% - Talk length is about 20min.
% - Style is ornate.



% Copyright 2004 by Till Tantau <tantau@users.sourceforge.net>.
%
% In principle, this file can be redistributed and/or modified under
% the terms of the GNU Public License, version 2.
%
% However, this file is supposed to be a template to be modified
% for your own needs. For this reason, if you use this file as a
% template and not specifically distribute it as part of a another
% package/program, I grant the extra permission to freely copy and
% modify this file as you see fit and even to delete this copyright
% notice. 


\mode<presentation>
{
  \usetheme{Warsaw}
  % or ...

  \setbeamercovered{transparent}
  % or whatever (possibly just delete it)
}

\usepackage{listings}
\usepackage[italian]{babel}
% or whatever

\usepackage[latin1]{inputenc}
% or whatever

\usepackage{times}
\usepackage[T1]{fontenc}
% Or whatever. Note that the encoding and the font should match. If T1
% does not look nice, try deleting the line with the fontenc.
\lstset{
  language=XML,
  basicstyle=\footnotesize,%\ttfamily,
  columns=fullflexible,
  morekeywords={encoding,
    xs:schema,xs:element,xs:complexType,xs:sequence,xs:attribute}
}

\title[Analisi di reti
metaboliche] % (optional, use only with long paper titles)
{Analisi di reti metaboliche basata su propriet\`a di connessione}

% \subtitle
% {Include Only If Paper Has a Subtitle}

\author[Massimo Nocentini] % (optional, use only with lots of authors)
{Massimo~Nocentini\\\texttt{massimo.nocentini@gmail.com}}
% - Give the names in the same order as the appear in the paper.
% - Use the \inst{?} command only if the authors have different
%   affiliation.

 \institute[UniversitaStudiFirenze] % (optional, but mostly needed)
 { Universit\`a degli Studi di Firenze }
%   \inst{1}%
%   Department of Computer Science\\
%   University of Somewhere
%   \and
%   \inst{2}%
%   Department of Theoretical Philosophy\\
%   University of Elsewhere}
% - Use the \inst command only if there are several affiliations.
% - Keep it simple, no one is interested in your street address.

\date[Tesi20120222] % (optional, should be abbreviation of conference name)
{Firenze, \today}
% - Either use conference name or its abbreviation.
% - Not really informative to the audience, more for people (including
%   yourself) who are reading the slides online

\subject{Theoretical Computer Science}
% This is only inserted into the PDF information catalog. Can be left
% out. 



% If you have a file called "university-logo-filename.xxx", where xxx
% is a graphic format that can be processed by latex or pdflatex,
% resp., then you can add a logo as follows:

\pgfdeclareimage[height=1.5cm]{university-logo}{logo/unifi}
\logo{\pgfuseimage{university-logo}}

% Delete this, if you do not want the table of contents to pop up at
% the beginning of each subsection:
\AtBeginSubsection[]
{
  \begin{frame}<beamer>{Outline}
    \tableofcontents[currentsection,currentsubsection]
  \end{frame}
}


% If you wish to uncover everything in a step-wise fashion, uncomment
% the following command: 

%\beamerdefaultoverlayspecification{<+->}


\begin{document}

\begin{frame}[plain]
  \titlepage
  % \begin{center}
  %   \includegraphics[scale=.065]{logo/unifi}
  % \end{center}
\end{frame}

\begin{frame}{Outline}
  \tableofcontents[pausesections]
  % You might wish to add the option [pausesections]
\end{frame}


% Structuring a talk is a difficult task and the following structure
% may not be suitable. Here are some rules that apply for this
% solution: 

% - Exactly two or three sections (other than the summary).
% - At *most* three subsections per section.
% - Talk about 30s to 2min per frame. So there should be between about
%   15 and 30 frames, all told.

% - A conference audience is likely to know very little of what you
%   are going to talk about. So *simplify*!
% - In a 20min talk, getting the main ideas across is hard
%   enough. Leave out details, even if it means being less precise than
%   you think necessary.
% - If you omit details that are vital to the proof/implementation,
%   just say so once. Everybody will be happy with that.

\section{Motivazioni}

\subsection{Analisi di reti metaboliche}

\begin{frame}{Definizione di rete metabolica}
  % - A title should summarize the slide in an understandable fashion
  %   for anyone how does not follow everything on the slide itself.
  \begin{definition}
    Una rete metabolica \`e un insieme di reazioni chimiche che si
    verificano all'interno di una cellula.
  \end{definition}
  \`E interessante studiarle per:
  \begin{itemize}
  \item capire le propriet\`a fisiologiche e biochimiche delle
    cellule;
  \item ricostruire le reazioni che avvengono all'interno di
    organismi, sia batteri che esseri umani;
  \item studiare il comportamento della cellula in relazione al
    contesto che la ospita.
  \end{itemize}
\end{frame}

\begin{frame}{Codifica delle reti}{usando il linguaggio \emph{SBML}}
  \textbf{SBML} (\textbf{S}ystems \textbf{B}iology \textbf{M}arkup
  \textbf{L}anguage) \`e un linguaggio orientato alla descrizione di
  sistemi in cui entit\`a biologiche sono oggetto di manipolazioni
  eseguite da processi fisici.
  \begin{figure}
    \includegraphics[scale=.6]{images/sbml-code-chunk.eps}
    \caption{Estratto della codifica del batterio Bartonella}
  \end{figure}
\end{frame}
  % You can create overlays\dots
  % \begin{itemize}
  % \item using the \texttt{pause} command:
  %   \begin{itemize}
  %   \item
  %     First item.
  %     \pause
  %   \item    
  %     Second item.
  %   \end{itemize}
  % \item
  %   using overlay specifications:
  %   \begin{itemize}
  %   \item<3->
  %     First item.
  %   \item<4->
  %     Second item.
  %   \end{itemize}
  % \item
  %   using the general \texttt{uncover} command:
  %   \begin{itemize}
  %     \uncover<5->{\item
  %       First item.}
  %     \uncover<6->{\item
  %       Second item.}
  %   \end{itemize}
  % \end{itemize}



\subsection{Previous Work}

\begin{frame}{Make Titles Informative.}
\end{frame}

\begin{frame}{Make Titles Informative.}
\end{frame}



\section{Our Results/Contribution}

\subsection{Main Results}

\begin{frame}{Make Titles Informative.}
\end{frame}

\begin{frame}{Make Titles Informative.}
\end{frame}

\begin{frame}{Make Titles Informative.}
\end{frame}


\subsection{Basic Ideas for Proofs/Implementation}

\begin{frame}{Make Titles Informative.}
\end{frame}

\begin{frame}{Make Titles Informative.}
\end{frame}

\begin{frame}{Make Titles Informative.}
\end{frame}



\section*{Summary}

\begin{frame}{Summary}

  % Keep the summary *very short*.
  \begin{itemize}
  \item
    The \alert{first main message} of your talk in one or two lines.
  \item
    The \alert{second main message} of your talk in one or two lines.
  \item
    Perhaps a \alert{third message}, but not more than that.
  \end{itemize}
  
  % The following outlook is optional.
  \vskip0pt plus.5fill
  \begin{itemize}
  \item
    Outlook
    \begin{itemize}
    \item
      Something you haven't solved.
    \item
      Something else you haven't solved.
    \end{itemize}
  \end{itemize}
\end{frame}



% All of the following is optional and typically not needed. 
\appendix
\section<presentation>*{\appendixname}
\subsection<presentation>*{For Further Reading}

\begin{frame}[allowframebreaks]
  \frametitle<presentation>{For Further Reading}
    
  \begin{thebibliography}{10}
    
  \beamertemplatebookbibitems
  % Start with overview books.

  \bibitem{Author1990}
    A.~Author.
    \newblock {\em Handbook of Everything}.
    \newblock Some Press, 1990.
 
    
  \beamertemplatearticlebibitems
  % Followed by interesting articles. Keep the list short. 

  \bibitem{Someone2000}
    S.~Someone.
    \newblock On this and that.
    \newblock {\em Journal of This and That}, 2(1):50--100,
    2000.
  \end{thebibliography}
\end{frame}

\end{document}

